%#!lualatex spotxcolor
\documentclass[luatex]{article}
\usepackage{shortvrb}\MakeShortVerb{\|}
\usepackage{hyperref}
\usepackage{xcolor}
\usepackage{spotxcolor}

\definespotcolor{DIC161}{DIC 161s*}{0, 0.64, 1, 0}

\title{\textsf{spotxcolor} Package}
\author{Munehiro Yamamoto}
\date{2026/03/14 v1.0}

\begin{document}
\maketitle

\begin{abstract}
This package provides robust spot color (e.g., DIC, PANTONE) support for the \textsf{xcolor} package across all major \TeX\ engines. It resolves structural PDF issues found in legacy packages and provides explicit fallback mechanisms for various drivers.
\end{abstract}

\section{Motivation \& Features}
While the legacy \textsf{spotcolor} package has been used for years, it suffers from structural PDF compatibility issues with modern \textsf{expl3} and \textsf{xcolor} updates. The \textsf{colorspace} package is an excellent modern alternative, but it lacks full support for the \texttt{dvipdfmx} driver.

To address these gaps, the \textsf{spotxcolor} package natively hooks into the \textsf{xcolor} package to provide a universal, print-safe solution. Its core features include:

\begin{itemize}
\item \textbf{Native \textsf{xcolor} Integration:} Use spot colors exactly like standard colors (e.g., \verb|\textcolor|, \verb|\pagecolor|), generating perfect PDF structures for \texttt{pdftex} and \texttt{luatex}.
\item \textbf{Universal Engine Support \& Safe Fallbacks:} Ensures safe CMYK fallback for \texttt{dvipdfmx} and \texttt{xetex} when using standard macros, while providing explicit injection commands (\verb|\SpotColor|) for true spot color output in \texttt{dvipdfmx}.
\item \textbf{Advanced Graphics (\textsf{TikZ}/\textsf{PGF}) Support:} Internally patches \textsf{PGF} to ensure print safety. It fully supports fill/stroke separation, fadings, shadings, blend modes, and safely forces uncolored patterns (including \texttt{patterns.meta}) to CMYK instead of hardcoded RGB.
\item \textbf{Seamless Decorations:} Works flawlessly with \textsf{colortbl} (zebra-striped tables) and \textsf{tcolorbox} (frames, backgrounds, shadows).
\item \textbf{Complete Backward Compatibility:} Perfectly emulates user commands from both \textsf{spotcolor} and \textsf{colorspace} packages.
\end{itemize}

\section{Requirements}
This package requires a modern \TeX\ Live environment (with an up-to-date \textsf{expl3}/\textsf{l3kernel}), alongside the \textsf{xcolor} and \textsf{iftex} packages.

\section{Loading the \textsf{spotxcolor} Package}

Load the \textsf{spotxcolor} package in your preamble. It automatically detects your engine via the \textsf{iftex} package.
\begin{quote}
\begin{verbatim}
\usepackage{spotxcolor}
\end{verbatim}
\end{quote}

\section{Usage}

\subsection{Defining and Using Spot Colors}
Use \verb|\definespotcolor| to register a spot color. The arguments are: \LaTeX\ name, PDF name, and CMYK values.
Once registered, you can use standard \textsf{xcolor} commands:

\begin{quote}
\begin{verbatim}
\definespotcolor{DIC161}{DIC 161s*}{0, 0.64, 1, 0}

\textcolor{DIC161}{This text is DIC 161.}
\end{verbatim}
\end{quote}

\subsection{Dvipdfmx Explicit Spot Colors}
Because the \texttt{dvipdfmx} backend strictly falls back to a CMYK representation when using standard \verb|\textcolor|, you can force true spot color output (creating an independent separation plate in the PDF) using the \verb|\SpotColor| command:

\begin{quote}
\begin{verbatim}
\SpotColor{DIC161}{1.0}
This text is 100\% DIC 161.

\SpotColor{DIC161}{0.5}
This text is 50\% DIC 161.
\end{verbatim}
\end{quote}

\subsection{Note for \texttt{(u)p\LaTeX} Users}
If you compile your document with \texttt{platex} or \texttt{uplatex} and generate a PDF via \texttt{dvipdfmx}, you \textbf{must} specify the \texttt{dvipdfmx} driver option globally in your document class:

\begin{quote}
\begin{verbatim}
\documentclass[dvipdfmx]{article}
\end{verbatim}
\end{quote}

Without this option, modern \textsf{expl3} and \textsf{xcolor} will default to the \texttt{dvips} driver, and the PDF objects for spot colors will not be generated correctly.

\subsection{Backward Compatibility}
You can safely reuse your legacy code and dictionaries written for the \textsf{spotcolor} package. They are automatically mapped to modern \textsf{spotxcolor} implementations:

\begin{quote}
\begin{verbatim}
\NewSpotColorSpace{DIC}
\AddSpotColor{DIC}{DIC161}{DIC\SpotSpace 161s*}{0 0.64 1 0}
\SetPageColorSpace{DIC}
\end{verbatim}
\end{quote}

\end{document}
